\section{共享年表之二：大国均势与内部裂解}

城门、渡口和盟坛并不属于某一位霸主。前七世纪末以后，晋、楚可以在一场会战中压倒对手，却不能把胜利变成对中原诸侯的常设管辖。小国在两强之间选择，大国也要为道路、军粮和盟友的犹疑付费；外部战事稍缓，国内掌握军政资源的卿大夫便有了更长的经营时间。这一百五十年所呈现的，正是这种秩序的反复缝合与逐渐松脱。

\begin{event}{公元前625年}{垂陇之盟}{经年可靠，叙事细节有争议}{垂陇}
\eventanchor{SA-0625-CHUILONG}{春秋·垂陇之盟}
晋的士縠与宋、陈、郑、鲁在垂陇会盟。卫事尚未平息，秦晋战后的关系也需要重新安置；对晋而言，把相邻诸侯请到同一处，并不只是礼仪上的点名，而是在黄河与中原道路之间确认谁愿意承担下一轮风险。对宋、陈、郑、鲁而言，会盟同样是提出条件、探问晋国保护是否仍可兑现的场合。

《春秋·文公二年》留下参盟者与年次的骨架，《左传》才把卫事和秦晋余波编入较完整的局势。盟辞、各国实际出了多少兵粮、此后是否都服从晋策，材料都没有保存。垂陇没有制造一个整齐的晋方集团；它只暂时把若干本来各有顾虑的国家织进同一张联络网，日后这张网仍须由救援、征伐和新的会盟反复修补。
\end{event}

\begin{event}{公元前597年}{邲之战}{可靠纪年}{邲}
\eventanchor{SA-0597-BI}{春秋·邲之战}
郑国刚在楚军围城的压力下转向，晋军便南下试图挽回盟友。邲地的会战因而不是孤立的两军相撞：楚要证明自己能保护新近归附者，晋则不能让中原诸侯相信求援无用。楚庄王击败晋军，战场上的退却很快变成各国重新计算道路、城邑和从属关系的消息。

《春秋·宣公十二年》只简记战争，《左传·宣公十二年》则把晋军内部的进退分歧、将帅责任和楚庄王的处置组织成富有伦理张力的叙事；《史记·楚世家》又从汉代的列国兴亡将其写入楚庄王的盛期。由这些材料可以确认败绩及其政治震动，却不能复原兵数、军令和人物说辞。直接后果是楚的北方声势上升；中期看，晋仍保有北方腹地和动员能力，双方转入更昂贵也更难一锤定音的均势竞争。
\end{event}

\begin{event}{公元前589年}{蜀地会盟}{经年可靠，叙事细节有争议}{蜀}
\eventanchor{SA-0589-SHU-ALLIANCE}{春秋·蜀地会盟}
齐、晋在鞌之战后重新排列彼此的位置，鲁却与楚、秦以及曹、邾、薛、鄫等国在蜀会盟。盟坛上并列的名字提醒人们，春秋外交并非只有两大阵营：刚经历战事的国家不愿把自己的道路和安全完全押给一方，小国尤其会借多重关系分散风险。

《春秋·成公二年》可以确认会盟的年次和参与者，《左传》提供相关叙事；“蜀盟改变了所有国家的战略”则只是有待检验的解释。盟约文字和兵役义务都已不存，不能从参盟名单推出永久的政治一致。不过，它使秦、楚和中原小国能够在晋齐冲突之后另开一条协商渠道，也说明霸权的边界常由未被吞并的小国划出。
\end{event}

\begin{event}{公元前568年}{晋侯会诸侯救陈}{经年可靠}{陈地通道}
\eventanchor{SA-0568-JIN-RESCUE-CHEN}{春秋·晋侯救陈}
楚伐陈，晋侯会合诸侯而救援。陈处在楚北进、晋南援都须经过的地带，故而一座城邑是否得到援军，关系到整条联盟的可信度。晋方把救陈说成盟友义务，也等于把诸侯的出兵、转运和等待，转化为一笔共同承担的成本。

《春秋·襄公五年》可据以系年，《左传》补出战争的叙事次序；援军究竟何时抵达、楚军因何撤退，不能由经文细化。此事没有使陈从此安稳，却令诸侯看见，\eventref{SA-0597-BI}{春秋·邲之战}以后晋仍能围绕一项救援调动网络。保护承诺一旦要靠反复征发来兑现，也就为日后的疲敝埋下了条件。
\end{event}

\begin{event}{公元前565年}{邢丘之会}{经年可靠，礼制含义有争议}{邢丘}
\eventanchor{SA-0565-XINGQIU}{春秋·邢丘之会}
诸侯会于邢丘，传世叙事将这次集会同朝聘次序相连。排列先后看似细小，实际牵涉谁须来见、谁可要求他人到场，以及晋能否把对郑、蔡等处的军事压力改写成可重复的程序。会盟若只靠临时出兵，成本会很快堆积；若能把出入、朝聘和名位做成惯例，霸主便似乎不必每次都拔剑。

然而《春秋·襄公八年》并未保存一部完整的朝聘章程，《左传》的礼制说明也不能直接当作固定官僚规章。邢丘更适宜被理解为一次秩序试验：晋试图以程序维持网络，诸侯则在同一程序中衡量自身负担。礼仪在这里不是战争的反面，而是为战争、征发与服从寻找较低成本的形式。
\end{event}

\begin{event}{公元前564年}{伐郑与戏盟}{经年可靠}{戏}
\eventanchor{SA-0564-XI-ALLIANCE}{春秋·戏盟}
郑向楚倾斜，晋率宋、卫、曹、莒、邾、滕、薛、杞、齐等国伐郑，继而在戏会盟。集结如此多的名字，不表示每一国都怀着同一目标；有些国家需要晋的保护，有些担心郑楚关系改变通道，也有些只是难以承受公开缺席的代价。

《春秋·襄公九年》保存了征伐与会盟的年次，《左传》使这场对郑的压力具有连续的行动线；盟辞、军队数额与诸侯实际履约程度却无从复原。相较于\eventref{SA-0565-XINGQIU}{春秋·邢丘之会}所追求的程序化秩序，戏盟露出其强制的一面：联盟能够惩罚摇摆者，却也把小国的人力和粮秣卷入大国纪律。
\end{event}

\begin{event}{公元前562年}{亳城北之盟}{经年可靠}{亳城北}
\eventanchor{SA-0562-BOCHENG-ALLIANCE}{春秋·亳城北之盟}
继续伐郑之前，晋与宋、卫、曹、齐等国盟于亳城北。军队要走多远、何时停下，并不是主盟国一声令下便能决定的事；盟会给予继续协同行军一个看得见的名义，也让各国能在行动中确认自己被期待承担何种位置。

《春秋·襄公十一年》与《左传·襄公十一年》使这次会盟可以系年，却未留下违约惩罚和补给分配的条款。因此不能把亳城北当作一份保存完好的军事同盟契约。它至少表明，\eventref{SA-0564-XI-ALLIANCE}{春秋·戏盟}后的晋方攻郑，仍须用新的盟誓延长合作；霸权不是一张签过便永久有效的盟书。
\end{event}

\begin{event}{公元前557年}{溴梁之盟}{经年可靠，参会层级有争议}{溴梁}
\eventanchor{SA-0557-MULIANG-ALLIANCE}{春秋·溴梁之盟}
晋悼公死后，平公面对的第一项难题不是抽象的“继承”，而是如何让已在外事与军政中占有位置的诸侯、卿大夫继续接受新的召集中心。诸侯会于溴梁，并由大夫盟，正发生在这种重新接续的时刻。国君是否亲临、谁代表谁歃血，都会影响承诺由谁背负。

《春秋·襄公十六年》记下会盟，《左传》提供更丰满的背景；二者都不能给出一张与会层级的完整名册，更不能据此断言大夫已取代国君。较可把握的是，外交不再只是一国君主彼此相见的事务，掌握军政和使节网络的人在盟会中愈发可见。它使晋方秩序得以续接，也使可累积为家门资源的职位更加重要。
\end{event}

\begin{event}{公元前554年}{祝柯之盟}{经年可靠}{祝柯}
\eventanchor{SA-0554-ZHUKE-ALLIANCE}{春秋·祝柯之盟}
晋围齐以后，与宋、卫、郑、曹、莒、邾、滕、薛、杞等盟于祝柯。围城与会盟相接，是要把一次军事行动变成更可持久的号召；对参盟者而言，则是以出席换取安全、发言权或避免成为下一个受压者的机会。

《春秋·襄公十九年》可以确认这场盟会，《左传》保存其叙事框架；盟辞没有传下，齐国内部的变化也不能由祝柯一役解释。祝柯的短期作用，是让晋继续显示可以召集诸侯；其限度同样显明：外部的号召并未阻止齐国内部权力重新分配，盟会无法替任何一国解决本国的继承和卿族问题。
\end{event}

\begin{event}{公元前553年}{澶渊之盟}{经年可靠}{澶渊}
\eventanchor{SA-0553-CHANYUAN}{春秋·澶渊之盟}
齐晋战后，晋、齐以及宋、卫、郑、曹等国盟于澶渊。河畔的盟坛并未抹去前一年的围困与猜忌，却提供了一种不必立刻继续用兵的联络框架。对夹在大国之间的国家而言，暂停公开冲突本身就是可以争取的收益。

《春秋·襄公二十年》与《左传·襄公二十年》能够证明会盟的发生，不能复原盟约内容，亦不能保证各方从此不再调整选择。澶渊应被看作一次可撤回的降温，而非永久和平。正因和平尚须不断续约，它才通向稍后更具规模、也更明确承认双方网络的\eventref{SA-0546-MIBING}{春秋·弭兵之会}。
\end{event}

\begin{event}{公元前546年}{宋向戌调停的弭兵之会}{可靠纪年}{宋}
\eventanchor{SA-0546-MIBING}{春秋·弭兵之会}
宋大夫向戌往来于晋、楚之间，终于使两强及若干诸侯在宋会盟。长期战争耗费的不是两个君主的声望而已：军队要过境，盟友要供给，郑、陈等居间国家还要不断承受改从的危险。弭兵的近因是这种消耗造成的僵局；使能条件则是宋仍能被各方当作中介，而非战场。

《春秋·襄公二十七年》给出会盟年次，《左传》保存向戌调停及较多安排；齐、秦并未完全纳入，盟辞和各国负担也不能逐条复原。因此所谓“弭兵”不是现代意义的裁军或集体安全制度，而是暂时承认晋、楚各有从属网络。它降低了部分直接战争，代价是各国内部掌握军政资源者有更多时间经营职位、封邑与随从；外部均势并没有消除内部裂解。
\end{event}

\begin{event}{公元前506年}{召陵会师侵楚}{经年可靠}{召陵}
\eventanchor{SA-0506-ZHAOLING-INVASION}{春秋·召陵侵楚}
蔡与楚的关系破裂后，晋、鲁、宋、蔡、卫、陈、郑等会于召陵并侵楚。它不能和同年的柏举混成一场战役：召陵是晋方重启的北方联盟行动，吴军的进攻则来自楚东面的另一条道路。两股压力在同一年出现，才使楚不得不同时计算多处通道和盟友。

《春秋·定公四年》与《左传·定公四年》可确证会盟、侵楚及其大致脉络，联军由谁统一指挥、各国出了多少兵却不见于材料。召陵的直接作用不是征服楚国，而是增加楚方的多线防务负担，并为蔡与吴的合作留下空间；这一条件随即在\eventref{SA-0506-BOJU}{春秋·柏举之战}中获得更剧烈的结果。
\end{event}

\begin{event}{公元前506年}{柏举之战，吴军攻入郢都}{可靠纪年}{柏举至郢都}
\eventanchor{SA-0506-BOJU}{春秋·柏举之战}
吴与蔡、唐联军由东面进入楚境，自小别、大别一带向柏举推进。楚令尹囊瓦在败势中离军，楚昭王离开郢都；吴军随后进入这座江汉政治中心。对于从水网与江淮通道崛起的吴国而言，把军队送进郢都，是进入天下竞争最醒目的证明；对于楚而言，都城失守却不是全部腹地和地方关系一同消失。

《春秋·定公四年》提供最小的编年骨架，《左传》是战役次序和楚廷危机的主要叙事来源，《史记·吴太伯世家》则以更晚的兴亡线加重人物色彩。三者都不足以精确核算路线、兵数或每一次命令。楚昭王得以依靠腹地与外援逐步恢复，吴军又难长期供给远离本土的占领，故而柏举显示的是远征可击穿大国，并不等于可以把大国编入自己的治理。孙武在此役的角色尤其须谨慎：《左传》的核心叙事未提孙子，后出的《史记》才使其居于显著位置；银雀山汉简可说明《孙子兵法》的早期文本传统，不能证明孙武必在柏举现场统军。
\end{event}

\begin{sourcenote}[从经年到战场：这些材料怎样相互校读]
此段会盟、征伐与战争，首先以《春秋》保存的鲁国年次、行动者和极简结果为脊柱；它不能替我们补出盟辞、兵数、军粮或动机。《左传》使垂陇、邲、弭兵、召陵与柏举成为可以连续阅读的政治叙事，却仍是后出的叙事文本，不是盟书、军令簿或会议速记。《国语》可用来观察国别说辞如何塑造合法性，《史记》可检验汉代怎样把列国关系整理为兴亡长线；二者都不应倒灌为春秋现场的逐字证词。历史地理与区域考古能够追问陈地、江淮和江汉道路怎样限制动员，却不能凭地形推定一场战役唯一的行军方案。
\end{sourcenote}

\begin{event}{公元前481年}{《春秋》记事终点}{传统分期}{鲁}
\eventanchor{SA-0481-END}{春秋·记事终点}
《春秋·哀公十四年》记“西狩获麟”，传统上以此年作为经文终点。经文在这里停止，并不表示诸侯间的战争、会盟、卿族竞争或地方征发同时停止；相反，吴越的战争仍在继续，晋国内的资源和城邑也仍在重新组合。

“获麟”在《公羊传》等经学传统中承载很强的象征意义。就材料史而言，前481年首先是传世《春秋》编年的终点，孔子及弟子传统、鲁国档案的存佚和经文为何止于此年，都是不同的问题。它适合标示一部文本的边界，不适合作为一条天然的历史断崖：会盟秩序留下的动员、从属与卿族问题，将在战国的国家重组中继续发作。
\end{event}

\begin{turningpoint}[制度得失：冻结冲突，不能冻结资源]
\term{所得}：从垂陇到弭兵，盟会使诸侯能协调救援、通行和停战，不必在每一次争端中立即投入全面战争；宋一类居间国家也因此拥有可被承认的外交位置。

\term{代价}：盟约缺少常设执行机构，仍要依赖国君、卿大夫、城邑和盟友不断给付。战争暂缓并不使兵役、封邑和军将职位回到公室，反而为它们在家门网络中累积提供时间。

\term{后续压力}：柏举显示远征与占领之间的距离；晋的卿族竞争则显示会盟无法收回国内资源。战国各国将以县邑、编户、军功和官吏来回答这些问题，但那不是从前481年突然开始的新世界。
\end{turningpoint}

\subsection{年表没有停止的地方：经文终点之外的持续竞争}

把前481年当作春秋的自然终点，容易让读者误以为获麟以后诸侯马上改用另一套政治语言。\eventref{SA-0481-END}{春秋·记事终点}只说明鲁国这部传世经文在此结束；同一段时间里，晋的卿族仍在出兵，齐的公室与田氏仍在重排位置，吴越的战事也没有因一部文本止笔而停下。年表的终点属于材料史，不是历史本身的一道城门。把两者分开，才能理解为什么战国的国家化既继承春秋的城邑、家门和动员，又必须寻找更能固定资源的办法。

晋的晚期行动尤其显示这种延续。\eventref{SA-0497-JIN-FACTION}{春秋·范氏中行氏赵氏之争}后，赵鞅、荀寅与士吉射所依赖的已不仅是朝廷中一时的宠信，而是晋阳、朝歌等城邑及其能够动员的关系。\eventref{SA-0493-TIE-BATTLE}{春秋·赵鞅与郑战于铁}、\eventref{SA-0490-ZHAO-ATTACK-WEI}{春秋·赵鞅伐卫}和\eventref{SA-0485-ZHAO-ATTACK-QI}{春秋·赵鞅侵齐}的连续记录不证明赵氏拥有一台完整的常设国家机器，却说明卿族冲突仍能越过国内边界，转为同郑、卫、齐相连的军事与外交行动。后来的三家分晋并非从一个停止运动的春秋世界里突然出现，而是这些城邑、职位和部众已经越来越难由晋君共同收束。

齐与吴越的后期也应放在同一尺度内。\eventref{SA-0489-TIAN-QI-REGICIDE}{春秋·陈乞弑齐孺子荼}使田氏的军政位置变得尖锐可见，却没有让齐国立即换名；\eventref{SA-0484-AILING}{春秋·艾陵之战}使吴的北上声望抵达高点，也没有替吴守住南方。到\eventref{SA-0482-HUANGCHI}{春秋·黄池之会}，夫差仍在北方争取盟序，越的压力却已把安全问题带回本土。历史并不按单一方向从“会盟”走向“国家”，而是在远征、继嗣、城邑控制和资源调度的不同节奏中不断改变。

晚春秋的年表因而需要双重阅读：一方面，它把可以系年的行动逐项留下，使后人不至于把复杂过程全写成抽象趋势；另一方面，它的断裂、短记和鲁国视角也提醒我们，任何一条连续叙事都带着保存者的边界。读者由此进入战国时不应期待一个完全陌生的开端，而应带着春秋遗下的难题前行：谁能让命令越过家门与城邑，谁为战争、迁徙和征调付出代价，王室承认又能在多大程度上追上已经发生的资源重组。
